\DocumentMetadata{lang=es-DO,testphase={phase-III,math}}
\documentclass[11pt]{scrartcl}
\usepackage{icv}

% -----------------------------------------------------------------------
% Migrado de legacy/Acueductos y Alcantarillados/Proyecto/
% 02_terminos_referencia.tex -- especificación de proyecto real
% completa, sin recortar.
% -----------------------------------------------------------------------
\icvsetup{
  asignatura  = {ICV 442-T -- Acueductos y Alcantarillados},
  titulocorto = {Términos de referencia},
  titulo      = {Términos de referencia},
  subtitulo   = {Proyecto de curso: sistemas de saneamiento para Las Charcas, Azua},
  profesor    = {Dr. Ricardo Hernández Moreira},
  periodo     = \icvciclo{2026}{3},
  version     = {v1},
}

\begin{document}
\icvmaketitle

\begin{icvobjectives}
  \item Desarrollar un diseño académico de sistemas de agua potable,
    alcantarillado sanitario y drenaje pluvial para una comunidad real,
    con memoria técnica, planos y modelación hidráulica que sustenten
    las decisiones de trazado y dimensionamiento.
\end{icvobjectives}

\section{Propósito del proyecto}
El proyecto consiste en desarrollar, en equipos pequeños, un diseño
académico de sistemas de agua potable, alcantarillado sanitario y
drenaje pluvial para la comunidad de Las Charcas, provincia Azua. El
producto esperado es un diseño detallado con memoria técnica, planos y
modelación hidráulica suficiente para sustentar decisiones de trazado y
dimensionamiento.

\section{Competencias que se evalúan}
\begin{center}
\begin{tabular}{p{0.16\linewidth} p{0.68\linewidth}}
  \toprule
  \textbf{Competencia} & \textbf{Aplicación en el proyecto} \\
  \midrule
  EA3 & Obtención y uso de información básica: población, topografía,
    lluvia, normativa y brechas de datos. \\
  EA4 & Aplicación de fundamentos de hidráulica y justificación
    normativa en los tres sistemas. \\
  EA5 & Consideración de factibilidad operativa, mantenimiento y
    coherencia general del diseño. \\
  \bottomrule
\end{tabular}
\end{center}

\section{Alcance}

\subsection{Sistemas incluidos}
\begin{itemize}
  \item \textbf{Agua potable}: fuente, conducción, almacenamiento, red
    de distribución y accesorios principales.
  \item \textbf{Alcantarillado sanitario}: colectores, pozos de visita,
    descarga final y bombeo solo si la topografía lo exige.
  \item \textbf{Drenaje pluvial}: captación superficial, colectores,
    estructuras de disipación y punto de descarga.
\end{itemize}

\subsection{Nivel de detalle esperado}
\begin{itemize}
  \item Memorias de cálculo trazables.
  \item Planos legibles en escala adecuada.
  \item Especificaciones técnicas de materiales y criterios de diseño.
  \item Modelos en EPANET y SWMM con resultados verificables.
\end{itemize}

\section{Marco normativo}
\begin{enumerate}
  \item Reglamento y criterios técnicos de INAPA.
  \item Reglamento del MOPC cuando establezca una exigencia más
    restrictiva o pertinente.
  \item Referencias internacionales solo como apoyo o para completar
    vacíos, nunca para desplazar sin justificación el marco local.
\end{enumerate}

\begin{icvwarning}[Atención]
No basta citar una institución de forma genérica. Las decisiones de
diseño deben vincularse, cuando sea posible, con artículos, tablas o
criterios identificables de la norma consultada.
\end{icvwarning}

\section{Información disponible}
\begin{center}
\begin{tabular}{p{0.29\linewidth} p{0.18\linewidth} p{0.37\linewidth}}
  \toprule
  \textbf{Dato} & \textbf{Estado} & \textbf{Observación} \\
  \midrule
  Modelo digital de terreno & Disponible & Malla de \qty{8}{\meter}
    por \qty{8}{\meter} suministrada para el curso \\
  Datos poblacionales de la ONE & Disponible & Base para proyección
    poblacional \\
  Series de precipitación & Parcial & Pueden solicitarse a INDOMET
    por vía institucional \\
  Curvas IDF & Disponible & Apéndice del reglamento de diseño de
    INAPA \\
  Información sobre fuentes existentes & Incierta & Puede
    investigarse con INAPA \\
  Usos de suelo y densidades & Parcial & Requieren interpretación a
    partir de cartografía e imágenes \\
  \bottomrule
\end{tabular}
\end{center}

\section{Organización de equipos}
\begin{itemize}
  \item Los equipos estarán formados por 2 o 3 estudiantes.
  \item Todos los integrantes deben participar en los tres sistemas.
  \item Se recomienda asignar roles de coordinación, modelación,
    documentación y planos, con posibilidad de rotación.
  \item Cada equipo debe mantener una bitácora semanal breve y útil.
\end{itemize}

\section{Entregables}
\begin{center}
\begin{tabular}{p{0.33\linewidth} p{0.17\linewidth} p{0.30\linewidth}}
  \toprule
  \textbf{Entregable} & \textbf{Formato} & \textbf{Observaciones} \\
  \midrule
  Memoria de proyección poblacional & PDF + hoja de cálculo &
    Entregable independiente \\
  Memoria técnica integrada & PDF & Incluye agua potable, sanitario
    y pluvial \\
  Planos del proyecto & PDF & Planta, perfiles y detalles
    principales \\
  Modelos hidráulicos & Archivos nativos + resumen PDF & EPANET y
    SWMM \\
  Presentación oral & PDF exportado & El archivo fuente puede
    elaborarse en Canva o programa equivalente \\
  Anexos & PDF y hojas de cálculo & Patrones, bitácoras, tabla
    maestra y soportes \\
  \bottomrule
\end{tabular}
\end{center}

\begin{icvnote}[Convención de nombres para entregas]
\texttt{ICV442\_T[Equipo]\_[Sistema]\_[Entregable]\_[Apellidos]\_v[versión].pdf}

Ejemplos de uso:
\begin{itemize}
  \item \texttt{ICV442\_T01\_AguaPotable\_MemoriaCalc\_HernandezPerez\_v1.0.pdf}
  \item \texttt{ICV442\_T02\_Sanitario\_PlanoPerfil\_GomezRuiz\_v1.0.pdf}
  \item \texttt{ICV442\_T03\_PresentacionOral\_TorresLopez\_v1.0.pdf}
\end{itemize}
\end{icvnote}

\section{Coordinación con instituciones}
Si el equipo decide gestionar información adicional, puede hacerlo
mediante solicitudes formales de acceso a la información dirigidas a
las oficinas correspondientes. Esa gestión es recomendable, pero no
obligatoria. Si no se obtiene respuesta dentro del plazo del curso, el
equipo debe documentar el intento y continuar con supuestos técnicos
razonables y explícitos.

\section{Criterios generales de aceptación}
\begin{itemize}
  \item El trabajo debe ser coherente entre sistemas, no una suma de
    tres ejercicios desconectados.
  \item Los supuestos deben estar declarados y ser compatibles con la
    información disponible.
  \item La presentación documental debe permitir que un tercero
    entienda qué se hizo, con qué datos y bajo qué criterios.
\end{itemize}

\end{document}
